\chapter{Исследование предсказателя переходов}
\label{ch:branch}

\section{Параметры исследования}

Программа \texttt{branch\_test} заполняет массив условий \texttt{cond} и массив данных \texttt{data} из \texttt{N} элементов и выполняет тестовый цикл:

\begin{lstlisting}[style=code, language=C++]
for (size_t r = 0; r < repeats; r++) {
    for (size_t i = 0; i < n; i++) {
        if (cond[i]) {
            sum += data[i];
        }
    }
}
\end{lstlisting}

Режимы:
\begin{itemize}
	\item предсказуемый --- условие имеет почти постоянное поведение: истинно для первой половины массива и ложно для второй;
	\item непредсказуемый --- используется заранее сформированная случайная последовательность условий.
\end{itemize}

В обоих режимах условие истинно примерно для половины элементов, поэтому объём вычислений одинаков.

Аналоги \texttt{branches} и \texttt{branch-misses} (таблица~\ref{tab:events-branch}) --- \texttt{INST\_BRANCH} и \texttt{BRANCH\_MISPRED\_NONSPEC}.
Доля ошибочных предсказаний:
$$
\mathit{Branch\ miss\ rate} = \frac{\mathit{branch\ misses}}{\mathit{branches}}.
$$

Как и в главе~\ref{ch:stride}, из всех показателей вычитается запуск с $\mathit{repeats} = 0$.

\begin{commands}
g++ branch_test.cpp -O3 -std=c++17 -o branch_test
sudo python3 scripts/branch.py
\end{commands}

\includescript{branch.py}{Измерения для двух режимов условия}

\section{Измеренные показатели}

\begin{table}[H]
\caption{Предсказуемые и непредсказуемые переходы}
\label{tab:branch}
\small
\begin{tabularx}{\textwidth}{|l|C|C|C|C|C|}
	\hline
	Режим & \texttt{branches} & \texttt{branch-\allowbreak misses} & Branch miss rate, \% & Время, с & \texttt{cycles} \\
	\hline
	Предсказуемый & \texttt{118750413} & \texttt{801} & \texttt{0.00} & \texttt{0.030} & \texttt{106547396} \\
	\hline
	Непредсказуемый & \texttt{118762321} & \texttt{47855817} & \texttt{40.30} & \texttt{0.223} & \texttt{913378231} \\
	\hline
\end{tabularx}
\end{table}

\section{Вывод}

В восьмом этапе получены значения branches, branch-misses, branch miss rate и времени выполнения для предсказуемого и непредсказуемого режимов.
При одинаковом числе переходов в непредсказуемом режиме доля ошибочных предсказаний составляет 40,3\% против практически нулевой в предсказуемом, а число тактов увеличивается в 8.6 раза.
